基于未校准二维平面访问管与理想电容组成的1T1C模型，建立受约束自然语言输入、Geant4电离沉积到TCAD等效源传递、匹配中性参考及测试前预测冻结的评估流程。统一50~$\Omega$预充版本完成四栅长、两粒子、两个方向下64个独立事件和128个0/1配对响应；训练64响应，域内留出与外推各32响应。在5~mV主阈值下，128个响应中96项正确、20项判据误读、8项临界、4项数值不确定，失败和缺失各为0与0项。域内留出的存储0/1岭回归RMSE分别为0.088294/0.094644~fC（电荷）和0.798415/0.838043~mV（裕量损失）。尺寸外推的存储0/1岭回归RMSE分别为0.068049/0.067155~fC（电荷）和0.519434/0.504929~mV（裕量损失）。14项氧化层case与预定两目标的匹配边界细化通过连续门槛，分类跨界和未选近界状态仍保留数值不确定。工程失败、未知退出码与四项确定性重试全部保全，重试和扩展不增加独立事件数。结论限于固定硅入口、二维等效源、受控准静态初态及所选电磁输运配置；不推断持久位翻转、商业DRAM可靠性或辐射截面。
